% ============================================================================
% CHAPTER 5: IMPLEMENTATION DETAIL
% ============================================================================
\chapter{Implementation Detail}

This chapter presents the detailed implementation of each system component with code structure, design patterns, key algorithms, and integration mechanisms.

\section{Project Directory Structure}

The project follows a modular package-based architecture as shown below:

\begin{lstlisting}[language={}, caption={Project Directory Structure}, label={lst:structure}]
ai-research-assistant/
|-- app/
|   |-- streamlit_app.py        # Main frontend (781 lines)
|   |-- templates/
|       |-- custom_styles.html
|       |-- css_print.css
|-- agents/
|   |-- __init__.py
|   |-- research_agent.py       # Tournament research (432 lines)
|   |-- methodology_agent.py    # Methodology proposals (328 lines)
|   |-- visualizer_agent.py     # Diagram rendering (2373 lines)
|   |-- invoice_agent.py        # Report synthesis (581 lines)
|   |-- autoresearch_agent.py   # Karpathy loop (293 lines)
|   |-- prompt_helpers.py       # Template rendering (8 lines)
|-- orchestrator/
|   |-- utils.py                # Utilities (15 lines)
|   |-- prompts/
|       |-- research_agent_template.txt
|       |-- invoice_agent_template.txt
|       |-- master_orchestrator_prompt.txt
|-- autoresearch/               # Karpathy experiment module
|   |-- train.py, program.md, results.tsv
|-- pdf/
|   |-- pdf_generator.py        # MD-to-PDF conversion
|-- output/                     # Generated reports
|-- .streamlit/config.toml      # Server configuration
|-- .env                        # API keys (gitignored)
|-- auto_research_worker.py     # Headless batch runner
\end{lstlisting}

\section{Research Agent Implementation}

The Research Agent (\texttt{agents/research\_agent.py}, 432 lines) is the core intelligence module. Its implementation comprises four key subsystems:

\subsection{Model Resolution and Failover}

The agent maintains a role-to-model mapping that assigns optimal models to each research perspective:

\begin{lstlisting}[language=Python, caption={Model Configuration}, label={lst:model_config}]
ROLE_MODEL_MAP = {
    "historical": "llama-3.1-8b-instant",
    "state_of_the_art": "qwen/qwen3-32b",
    "ongoing_emerging": "llama-3.1-8b-instant",
}
FALLBACK_MODELS = ["llama-3.1-8b-instant", "qwen/qwen3-32b"]
\end{lstlisting}

The \texttt{call\_llm()} function implements multi-model failover with rate-limit detection:

\begin{lstlisting}[language=Python, caption={LLM Call with Failover (Simplified)}, label={lst:call_llm}]
def call_llm(prompt, role, temperature=0.3):
    primary = resolve_model(ROLE_MODEL_MAP.get(role))
    models = [primary] + [m for m in FALLBACK_MODELS
              if m != primary]
    for model in models:
        for attempt in range(2):
            try:
                completion = client.chat.completions.create(
                    model=model,
                    messages=[{"role": "system",
                               "content": SYSTEM_MSG},
                              {"role": "user",
                               "content": prompt}],
                    temperature=temperature,
                    max_tokens=8000, timeout=90)
                content = completion.choices[0].message.content
                if content and content.strip():
                    return content
            except Exception as e:
                if "429" in str(e):  # Rate limit
                    break  # Switch model immediately
    raise RuntimeError("All models failed")
\end{lstlisting}

\subsection{Balanced-Brace JSON Extraction}

The \texttt{\_balanced\_extract()} function (lines 176--236) implements a finite-state parser that tracks brace depth, string state, and escape sequences to extract the first complete JSON object from potentially malformed LLM output. When truncation is detected (depth $> 0$ at end of string), the function performs structural repair by:
\begin{enumerate}[label=\arabic*., leftmargin=2cm]
\item Removing trailing commas or colons
\item Closing unclosed strings with a quote character
\item Re-scanning the repaired string to count open braces/brackets
\item Appending closing characters in reverse stack order
\end{enumerate}

\subsection{Tournament Execution via ThreadPoolExecutor}

The \texttt{run\_research\_agent()} function spawns three parallel experiments using Python's \texttt{concurrent.futures.ThreadPoolExecutor}:

\begin{lstlisting}[language=Python, caption={Parallel Tournament Execution}, label={lst:tournament}]
with ThreadPoolExecutor(max_workers=3) as executor:
    futures = [executor.submit(_run_experiment, i)
               for i in range(3)]
    for future in as_completed(futures):
        result = future.result()
        if result["score"] > best_score:
            best_score = result["score"]
            best_result = result
\end{lstlisting}

\section{Visualizer Agent Implementation}

The Visualizer Agent (\texttt{agents/visualizer\_agent.py}, 2373 lines) is the largest module, implementing programmatic rendering of six diagram types using Pillow's \texttt{ImageDraw} API.

\subsection{Theme Engine}

The system defines 8 complete visual themes (Table~\ref{tab:themes}), each containing: background colors, card colors, border colors, accent colors, text colors, per-layer color mappings (client, gateway, service, AI/ML, data, infrastructure), step colors, background patterns (dots, grid, diagonal, hexagon, none), and layout styles (horizontal\_layers, vertical\_columns, radial, zigzag).

\begin{table}[H]
\centering
\caption{Visual Theme Catalog}
\label{tab:themes}
\begin{tabularx}{\textwidth}{|c|l|l|X|}
\hline
\textbf{\#} & \textbf{Theme Name} & \textbf{BG Pattern} & \textbf{Layout Style} \\
\hline
0 & Dark Cyber & Dots & Horizontal Layers \\
1 & Clean White & Grid & Horizontal Layers \\
2 & Midnight Ocean & Diagonal & Vertical Columns \\
3 & Warm Sunset & Hexagon & Zigzag \\
4 & Arctic Frost & None & Horizontal Layers \\
5 & Neon Matrix & Grid & Vertical Columns \\
6 & Coral Reef & Dots & Zigzag \\
7 & Deep Space & Hexagon & Horizontal Layers \\
\hline
\end{tabularx}
\end{table}

\subsection{Architecture Diagram Rendering}

The \texttt{\_render\_architecture()} function generates layered system diagrams with dynamically calculated dimensions based on content:

\begin{equation}
H = H_{\text{header}} + (N_{\text{layers}} \times H_{\text{layer}}) + ((N_{\text{layers}} - 1) \times G_{\text{gap}}) + H_{\text{footer}}
\label{eq:diagram_height}
\end{equation}

where $H_{\text{header}} = 240$px, $H_{\text{layer}} = 400$px, $G_{\text{gap}} = 60$px, $H_{\text{footer}} = 100$px, and $N_{\text{layers}}$ ranges from 4 to 6 depending on system category. The width is computed as:

\begin{equation}
W = \max(3200, W_{\text{sidebar}} + 100 + N_{\text{maxcomp}} \times (W_{\text{comp}} + G_{\text{comp}}) + 100)
\label{eq:diagram_width}
\end{equation}

\section{Streamlit Frontend Implementation}

The Streamlit frontend (\texttt{app/streamlit\_app.py}, 781 lines) implements a WebSocket-safe polling architecture. The critical design pattern uses \texttt{time.sleep(0.5)} polling loops instead of blocking \texttt{future.result(timeout=N)} calls:

\begin{lstlisting}[language=Python, caption={WebSocket-Safe Agent Execution}, label={lst:websocket}]
def _run_agent_safe(fn, *args, timeout_seconds=600):
    pool = ThreadPoolExecutor(max_workers=1)
    future = pool.submit(fn, *args)
    deadline = time.time() + timeout_seconds
    while not future.done():
        if time.time() > deadline:
            future.cancel()
            return None, "Timeout"
        time.sleep(0.5)  # Yield to Streamlit loop
    return future.result(timeout=0), None
\end{lstlisting}

This architecture prevents WebSocket ping/pong frame starvation that causes ``Connection error'' disconnects after 60--120 seconds of blocking calls.
